% T25 source: future_work.tex
For a fixed graph, the proposed model solves the forward problem: given the
work structure and the parameters of the human--agent cell, it estimates the
achievable speedup and constructs a schedule. The inverse problem is a natural
extension. If only part of the graph $G_t$ is known at time $t$, further
decomposition can be selected in light of the critical path, the required
parallelism, and the limited human-attention resource. The model would thereby
serve not only to evaluate a completed plan but also as a criterion for
constructing the as-yet unrevealed part of the project.

Such a policy can follow a rolling-horizon approach. After a result reveals
new work or dependencies, $G_t$, the set of ready tasks, and the phase
estimates are updated, and the unfinished part is rescheduled. Research tasks
can then be represented as exploratory vertices whose results reveal a set of
feasible continuations $\mathcal D_t$. For each $d\in\mathcal D_t$, the
following surrogate estimate can be used:

\[
\Phi_t(d)=
\max\left\{
\frac{W_{4,t}(d)+\widehat Q_t(d)}{P},
L_{4,t}(d),
\gamma(P)H_t(d)
\right\},
\qquad
d_t^*\in\operatorname*{arg\,min}_{d\in\mathcal D_t}\Phi_t(d),
\]

where $\widehat Q_t(d)$ estimates future waiting that is not captured by the
structural bounds already known. A decomposition variant must preserve the
original workload and acceptance criterion, while its additional task
specification, verification, and integration overhead must be included in
$W_{4,t}(d)$ and $H_t(d)$. If an empirically established manageable limit
$P_{\max}^{\mathrm{attn}}$ on the number of streams that can be supervised
concurrently is available, exceeding that limit should be prohibited,
and the proximity of useful parallelism to $P_{\max}^{\mathrm{attn}}$ should
be used as a secondary criterion after minimizing $\Phi_t$. Constructing
feasible graph transformations, estimating $\widehat Q_t(d)$, and deriving
guarantees for such an adaptive policy constitute a separate joint
decomposition-and-scheduling problem; in this article, the criterion is
presented only as a sketch of a future theory.

A separate direction is to extend the model to other multi-agent and swarm
architectures. The closest extension of the current formulation is a
hierarchical variant in which an orchestrator agent replaces the developer in
the inner loop: it decomposes the objective, assigns tasks to worker agents,
services their requests, and accepts their results. The structural analogy is
direct: the orchestrator is likewise a shared resource of finite capacity, and
requests from worker agents form a queue and may locally block their streams.
Human attention may then remain only at the boundaries of the overall
process---when setting the overall objective and performing final acceptance.

A different formulation arises under algorithmic orchestration with auditing
and self-checking. The macro-level task graph may remain deterministic, but
each vertex expands into a local loop of ``execution---audit---conditional
rework---repeated audit.'' If such loops are self-contained and do not use a
shared orchestrator, they do not block one another, and the internal human
resource of capacity 1 disappears. In the terms of the present model, audit
and rework become agent phases: $h_i$ decreases for the internal vertices,
while $a_i$ increases; the sign of the change in $k_i=a_i+h_i$ depends on the
balance between the human effort saved and the additional agent work. With a
fixed number of passes, such a loop can be unfolded into an ordinary phase
DAG; conditional termination and an unknown number of rework passes require a
stochastic model of rework.

A more general extension should describe an architecture by a graph of
resources and interactions rather than by the number of agents. Hierarchical
orchestrators, decentralized swarms, mutual auditing, specialized planners and
verifiers, shared memory, and heterogeneous resources create different
bottlenecks and different forms of coordination overhead. Comparing such
systems requires replacing the star-shaped ``one developer---$P$ agents''
scheme with a heterogeneous multilevel scheduling model and evaluating the
configuration as a whole: its topology, routing rules, audit resources, result
quality, and the cost of additional self-checking cycles. Multiple
human--agent cells require additional cross-cell resources and therefore
cannot be reduced to an increase in $P$ at M4.

Another direction is relative and interval calibration for organizations that
lack a history of precise measurements. Under cold-start conditions, the unit
$X$ can be interpreted not as an hour but as a notional baseline task, after
which $Z_i$, $k_i$, and $h_i$ can be elicited from experts relative to that
unit. If two configurations $A$ and $B$ are described on the same scale and all
of their agent and human-attention phases contain the same unknown common
factor $\kappa>0$, Proposition~\ref{prop:invariance} implies

\[
\frac{T_A^*(\kappa)}{T_B^*(\kappa)}=
\frac{T_A^*(1)}{T_B^*(1)}.
\]

Thus, even coarse quantitative relative estimates can be useful for selecting
a faster work-organization scheme without promising an exact completion time
in hours. Further research should replace point estimates with intervals for
$k_i$ and $h_i$, identify conditions under which one variant robustly
dominates another, and characterize cases in which uncertainty in the relative
profile can alter the ranking.

Another direction is a stochastic and robust formulation with rework,
failures, learning, and interval estimates of $k$, $h$, $C$, and $\gamma$. It
should provide not only the expected completion time but also quantiles, the
probability of missing a deadline, and conditions under which one
configuration robustly dominates another.

Finally, empirical calibration on traces of real software work is necessary.
It should measure autonomous intervals, human-attention intervals, waiting,
rework, and result quality separately in order to test the forms of $C(P)$ and
$\gamma(P)$ and distinguish effective parallelism from the nominal number of
running agents.
